% s-235-p235-m002; original PDF 235, 236, 237
\Aphorism{s-235-p235-m002}{329}
Ob wirklich in der Kunst die Persönlichkeit soviel bedeutet, wie die Künstler selbst glauben, nämlich Alles? Man kann dies leugnen, ohne darum den Massen, dem „öconomischen Process“ oder sonstwelchen kunstfremden Mächten einen Antheil an der Kunstentwicklung einzuräumen; diese beschränkte Alternative zwischen ideellen und materiellen Motoren ist genau so wenig vorhanden wie in der Naturwissenschaft die zwischen einem denkenden Schöpfer und brutalem Zufallsspiel. Muss man die Weltgeschichte für prästabilirte Harmonie halten, wenn man sie nicht als sinnlose Configuration ansehen will, wie es Schopenhauer verlangt? Und ist das Kunstwerk seiner mystischen Einzigkeit entkleidet, wenn man die persönliche Urheberschaft für unzureichend hält, es hervorzubringen? Gewiss nicht, antworten die Anbeter des „Genius“; nur ist das Überpersönliche, das den Künstler im gegebenen Falle befruchtet und befähigt hat, eben wieder ein eminent Persönliches, die Inspiration nur dieses Einen, die ihm gnadenweise, aber ausschliesslich gewordene Offenbarung und Berufung, solche zu verkünden — oder wie man sich, mehr oder minder metaphysisch, ausdrücken will. Und die Zeit? die Tradition? die Stetigkeit in der Stellung und Bewältigung neuer Aufgaben? An diese wirkenden Kräfte wird unter Künstlern selten gedacht; deren Antheil am Kunstwerke herauszusondern überlässt man der Ästhetik, die ihrerseits mit einer gewissen Rancune ausschliesslich an solche historische Bedingungen anknüpft und das Persönliche ganz fort haben möchte. Aber Beides lebt im Kunstwerk, Zeitwille und Einzelwille, und weder die Revolution ohne Beethoven, noch Beethoven ohne die Revolution hätte die Blüthe einer Eroica getrieben; Raffael war kein ausführender Handwerker und Papst Leo kein zahlender Auftraggeber, sondern Beide nahmen und gaben, und hinter Beiden stand der Geist der Renaissance, der ihrer so wenig entrathen konnte wie sie seiner. Der „Geist der Renaissance“? Da hätten wir ja die Mystik wieder, nur in anderem Gewande! Wer es mit diesen Zeitgeistern nicht hält, der erwäge doch, ob wirklich die Kunst nicht einer spontanen Fortpflanzung fähig sein sollte, die von selbst die schöpferischen Individuen, deren sie zum Sichtbarwerden bedarf, ins Leben ruft, ob nicht durch Accumulation gewisser Gedanken- und Gefühlstendenzen eine Energie der Weiterwirkung sich aufhäuft, die bis zu einem gewissen Grade vom Zufall genialer Einzelwesen unabhängig ist. Auf wissenschaftlichem Felde ist dies kaum zu bezweifeln. Die Entwicklung der neueren Astronomie zeigt ein Verhalten von der erstaunlichsten Zweckmässigkeit und Folgerichtigkeit; in ganz kurzem Zeitraum wird hintereinander die principielle Auffassung der himmlischen Bewegungen, ihre ziffermässige Festlegung, ihre geometrische und ihre mechanische Interpretation gewonnen, sodass kein Glied dieser Entdeckungenfolge fehlen durfte, ohne die nachfolgenden unmöglich zu machen. Ob diese verblüffende Teleologie wirklich an jenes geweihte Pentagramm, an das Zufallsspiel mit den fünf Glückswürfen Copernicus, Tycho, Kepler, Galilei, Newton geknüpft war? Oder ob wir nicht dieselbe Sache in persönlich anderer Färbung besässen, wenn in einem anderen Kopfe als dem des Frauenburger Canonicus der pythagoräische Funke gezündet hätte? Logische Zusammenhänge von dieser Schoss- und Triebkraft giebt es aber auch in der Kunst (man denke an die deutsche Musik von Bach bis Wagner): selbstthätig sich weiterspinnende Entwickelungsreihen und Gedankenfolgen, in die das schöpferische Einzelwesen mehr hineingezogen wird als activ eingreift. Duft, Farbe und Eigenart des Kunstwerks stammt vom Persönlichen, aber seinen historischen Rang, sein „Argument der Epoche“ bringt es von anderswoher mit, damit beschenkt es Den, welchem es seine Gestaltung verdanken will!
\par

